% Generated mechanically from the frozen fdl-zh.tex target.
% Do not edit this generated file; edit the bound locale source instead.

% OK, start here.
%
%---------------------------------------------------------------------
%\tableofcontents

\setcounter{chapter}{115}
\chapter{GNU 自由文档许可证（GNU Free Documentation License）}

\phantomsection
\label{fdl-zh-section-phantom}

\begin{center}

      1.2 版，2002 年 11 月


版权所有 \copyright 2000, 2001, 2002 Free Software Foundation, Inc.

\bigskip

    51 Franklin St, Fifth Floor, Boston, MA  02110-1301  USA

\bigskip

人人均可复制和分发本许可证文档的逐字副本，
但不得对其作任何更改。
\end{center}

This is an unofficial translation of the GNU Free Documentation License into
Simplified Chinese. It was not published by the Free Software Foundation, and
does not legally state the distribution terms for documentation that uses the
GNU FDL---only the original English text of the GNU FDL does that. However, we
hope that this translation will help Simplified Chinese speakers understand
the GNU FDL better.

这是 GNU 自由文档许可证的非官方简体中文译文。它并非由自由软件基金会发布，也不构成采用
GNU FDL 的文档所适用的法律分发条款；只有 GNU FDL 的英文原文具有该等法律效力。不过，
我们希望本译文能帮助简体中文读者理解 GNU FDL。本译文必须与英文原文一并提供；如译文与
英文原文有任何分歧，以英文原文为准。

You may publish this translation, modified or unmodified, only under the terms
at https://www.gnu.org/licenses/translations.html.

只有遵守 https://www.gnu.org/licenses/translations.html 所列条款，方可发布本译文，
无论是否经过修改。


\begin{center}
{\bf\large 前言（Preamble）}
\end{center}

本许可证旨在使手册、教科书或其他具有实际功能和用途的文档获得自由意义上的“自由”：
确保每个人都切实享有复制和再分发该文档的自由，无论是否修改，也无论出于商业还是非商业
目的。其次，本许可证为作者和出版者保留一种使其工作获得署名认可的方式，同时不使他们因
他人所作的修改而承担责任。

本许可证是一种“copyleft”（著佐权）许可证；这意味着文档的衍生作品本身也必须在同一
意义上保持自由。它与 GNU 通用公共许可证相配合；后者是为自由软件设计的著佐权许可证。

我们设计本许可证，是为了将其用于自由软件的手册，因为自由软件需要自由文档：自由程序
应当附有向读者提供与该软件相同自由的手册。不过，本许可证并不限于软件手册；任何文字作品
都可以采用它，而不论主题为何，也不论是否以印刷书籍出版。我们尤其建议以教学或参考为目的
的作品采用本许可证。


\section{适用范围与定义（APPLICABILITY AND DEFINITIONS）}
\label{fdl-zh-section-applicability-and-definitions}

本许可证适用于以任何媒介承载的任何手册或其他作品，只要其中载有版权持有人作出的声明，
说明该作品可依本许可证条款分发。该声明授予一项全球范围、免版税、期限不受限制的许可，
允许依本文所述条件使用该作品。下文所称的\textbf{“文档”}，指任何此类手册或作品。
任何公众成员都是被许可人，本文以\textbf{“您”}称呼被许可人。若您以版权法要求获得许可
的方式复制、修改或分发该作品，即表示您接受本许可证。

文档的\textbf{“修改版”}，指任何包含该文档或其一部分的作品，无论其内容是逐字复制，
还是经过修改和／或翻译成另一种语言。

\textbf{“次要章节”}，指文档中有标题的附录或前置章节；该章节只讨论文档的出版者或作者
与文档总体主题（或相关事项）之间的关系，不包含任何可直接属于该总体主题的内容。（因此，
若文档的一部分是数学教科书，次要章节不得讲解任何数学。）这种关系可以是与该主题或相关
事项的历史联系，也可以是对这些事项所持的法律、商业、哲学、伦理或政治立场。

\textbf{“不变章节”}，指某些次要章节，其标题在说明文档依本许可证发布的声明中被指定为
不变章节的标题。若某一章节不符合上述次要章节的定义，就不得将其指定为不变章节。文档可以
没有不变章节。若文档未指明任何不变章节，则不存在不变章节。

\textbf{“封面文字”}，指说明文档依本许可证发布的声明中列为前封面文字或后封面文字的
某些简短文字。前封面文字最多为 5 个词，后封面文字最多为 25 个词。

文档的\textbf{“透明”}副本，指采用一种机器可读格式表示的副本；该格式的规范可供公众
取得，并适合使用通用文本编辑器直接修订文档，或者（对于由像素组成的图像）使用通用绘图
程序，或者（对于线条图）使用某种广泛可得的绘图编辑器修订；该格式还适合输入排版程序，
或适合自动转换为多种可输入排版程序的格式。若某一原本属于透明格式的文件，其标记方式或
无标记方式被特意安排成阻碍或劝阻读者日后修改，则该副本并不透明。若图像格式被用于承载
大量文字，该格式并不透明。不是“透明”的副本称为\textbf{“不透明”}副本。

适合作为透明副本的格式包括：无标记的纯 ASCII、Texinfo 输入格式、LaTeX 输入格式、采用
公开可得 DTD 的 SGML 或 XML，以及符合标准、为便于人工修改而设计的简单 HTML、PostScript
或 PDF。透明图像格式包括 PNG、XCF 和 JPG。不透明格式包括：只能由专有文字处理器读取和
编辑的专有格式；其 DTD 和／或处理工具并非普遍可得的 SGML 或 XML；以及某些文字处理器
仅为输出目的而生成的 HTML、PostScript 或 PDF。

对于印刷书籍，\textbf{“标题页”}指标题页本身，以及为清晰容纳本许可证要求在标题页出现
的材料所必需的后续页面。对于本身没有标题页的格式，“标题页”指正文开始之前、接近作品标题
最醒目处的文字。

一个\textbf{“标题为 XYZ”}的章节，是指文档中有标题的子单元，其标题恰为 XYZ，或者在
另一种语言对 XYZ 的译文之后以括号包含 XYZ。（这里的 XYZ 代表下文提到的某个特定章节名，
例如\textbf{“Acknowledgements”}（致谢）、\textbf{“Dedications”}（献辞）、
\textbf{“Endorsements”}（背书）或\textbf{“History”}（历史）。）当您修改文档时，
\textbf{“保留该标题”}，是指该章节仍依本定义构成“标题为 XYZ”的章节。

文档可以在声明本许可证适用于该文档的通知旁列出担保免责声明。这些担保免责声明视为通过
引用纳入本许可证，但仅限于免除担保责任：这些担保免责声明可能产生的任何其他含义均属无效，
且不影响本许可证的含义。


\section{逐字复制（VERBATIM COPYING）}
\label{fdl-zh-section-verbatim-copying}

您可以在任何媒介中复制和分发文档，无论出于商业还是非商业目的，条件是：所有副本均重现
本许可证、版权声明以及说明本许可证适用于该文档的许可声明；并且您不得在本许可证的条件
之外附加任何其他条件。您不得使用技术措施妨碍或控制对您制作或分发之副本的阅读或进一步
复制。不过，您可以因提供副本而收取报酬。若您分发的副本数量足够大，还必须遵守第 3 节的
条件。

您也可以按上述相同条件出借副本，并可以公开展示副本。


\section{大量复制（COPYING IN QUANTITY）}
\label{fdl-zh-section-copying-in-quantity}


若您出版的文档印刷副本（或采用通常带有印刷封面的媒介制作的副本）超过 100 份，并且文档
的许可声明要求封面文字，则必须为副本加上清晰易读地载有全部封面文字的封面：前封面载有
前封面文字，后封面载有后封面文字。两个封面还必须清晰易读地表明您是这些副本的出版者。
前封面必须列出完整标题，标题中的所有词都应同样醒目可见。除上述内容外，您可以在封面上
增添其他材料。若改动仅限于封面，同时保留文档标题并满足这些条件，则在其他方面可视为逐字
复制。

若任一封面所要求的文字太长，无法清晰排入，您应将所列文字中的前若干项（在合理范围内尽量
多排）置于实际封面，并将其余内容续排在相邻页面。

若您出版或分发的文档不透明副本超过 100 份，则必须满足以下二者之一：随每一份不透明副本
附上一份机器可读的透明副本；或者在每一份不透明副本内或随附材料中说明一个计算机网络位置，
使一般网络公众能够通过公共标准网络协议下载文档的完整透明副本，且该透明副本不含附加材料。
若采用后一方案，则在开始大量分发不透明副本时，您必须采取合理审慎的措施，确保该透明副本
在所述位置持续可访问，直至您最后一次向公众分发该版本的不透明副本（无论直接分发还是通过
代理商或零售商分发）之后至少一年。

我们请求但不强制要求您在大量再分发副本之前及早联系文档作者，使其有机会向您提供文档的
更新版本。


\section{修改（MODIFICATIONS）}
\label{fdl-zh-section-modifications}

您可以按照上述第 2 节和第 3 节的条件复制和分发文档的修改版，但须将修改版完全依本许可证
发布，并由修改版承担“文档”的角色，从而允许任何持有修改版副本的人分发和修改该修改版。
此外，您必须在修改版中完成以下事项：

\begin{itemize}
\item[A.]
   在标题页（以及封面，如有）使用一个不同于文档标题和任何先前版本标题的标题；先前版本的
   标题（如有）应列在文档的“History”（历史）章节中。若某一先前版本的原出版者许可，您
   可以使用与该版本相同的标题。

\item[B.]
   在标题页上以作者身份列出一个或多个对修改版中的修改负有作者责任的个人或实体，同时列出
   文档至少五位主要作者（若主要作者不足五位，则列出全部主要作者），除非这些作者免除您的
   此项义务。

\item[C.]
   在标题页上以出版者身份写明修改版出版者的名称。

\item[D.]
   保留文档的全部版权声明。

\item[E.]
   在其他版权声明旁，为您所作的修改添加适当的版权声明。

\item[F.]
   紧接版权声明之后，按照下文附录所示格式加入许可声明，准许公众依本许可证条款使用修改版。

\item[G.]
   在该许可声明中保留文档许可声明所列不变章节和所要求封面文字的完整清单。

\item[H.]
   包含一份未经改动的本许可证副本。

\item[I.]
   保留标题为“History”（历史）的章节，保留其标题，并在其中加入一项，至少写明标题页所列
   修改版的标题、年份、新作者和出版者。若文档没有标题为“History”的章节，则创建该章节，
   写明文档标题页所列的标题、年份、作者和出版者，然后再按前一句所述加入一项描述修改版。

\item[J.]
   保留文档中为公众访问其透明副本而提供的网络位置（如有），并同样保留文档所依据的先前版本
   中所给的网络位置。这些信息可以放在“History”章节中。若某一作品在文档本身之前至少四年
   已经出版，或者该网络位置所指版本的原出版者允许，您可以省略该作品的网络位置。

\item[K.]
   对于任何标题为“Acknowledgements”（致谢）或“Dedications”（献辞）的章节，保留该章节
   的标题，并在该章节中保留其中每一项贡献者致谢和／或献辞的全部实质内容与语气。

\item[L.]
   保留文档的全部不变章节，不得改动其文字和标题。章节编号或同等标识不视为章节标题的一部分。

\item[M.]
   删除任何标题为“Endorsements”（背书）的章节。修改版不得包含此类章节。

\item[N.]
   不得将任何既有章节改名为“Endorsements”，也不得将其改成与任何不变章节标题冲突的标题。

\item[O.]
   保留所有担保免责声明。
\end{itemize}

若修改版包含新的前置章节或附录，且这些部分符合次要章节的定义并且不含从文档复制的材料，
您可以自行选择将其中一些或全部指定为不变章节。为此，请把它们的标题加入修改版许可声明中的
不变章节清单。这些标题必须与其他所有章节标题不同。

您可以加入标题为“Endorsements”（背书）的章节，但其中只能包含各方对您的修改版所作的
背书，例如同行评审声明，或说明某一组织已认可该文本为某项标准的权威定义。

您可以在修改版封面文字清单的末尾加入一段最多五个词的前封面文字，以及一段最多二十五个词
的后封面文字。任何一个实体（或由其安排）只能加入一段前封面文字和一段后封面文字。若文档
已经在同一封面包含先前由您加入、或由您所代表的同一实体安排加入的封面文字，您不得再加入
另一段；但在取得添加旧文字之先前出版者的明确许可后，可以替换旧文字。

文档的作者和出版者并未通过本许可证准许他人使用其姓名作宣传，也未准许他人声称或暗示其为
任何修改版背书。


\section{合并文档（COMBINING DOCUMENTS）}
\label{fdl-zh-section-combining-documents}


您可以依上述第 4 节为修改版规定的条款，将文档与其他依本许可证发布的文档合并，但条件是：
合并作品中必须包含所有原文档的全部不变章节，且不得改动；必须在合并作品的许可声明中将它们
全部列为不变章节；并且必须保留其全部担保免责声明。

合并作品只需包含本许可证的一份副本，多个完全相同的不变章节可以合并为一份。若存在标题相同
但内容不同的多个不变章节，应在每一章节标题末尾的括号内加入该章节原作者或原出版者的姓名
（如已知），否则加入一个唯一编号，使各标题彼此不同。合并作品许可声明中的不变章节清单也须
对相应章节标题作同样调整。

合并时，必须把各原文档中标题为“History”（历史）的章节合并成一个标题为“History”的章节；
同样合并所有标题为“Acknowledgements”（致谢）的章节和所有标题为“Dedications”（献辞）的
章节。必须删除所有标题为“Endorsements”（背书）的章节。

\section{文档合集（COLLECTIONS OF DOCUMENTS）}
\label{fdl-zh-section-collections-of-documents}

您可以制作由本文档和其他依本许可证发布的文档组成的合集，并以合集内的一份许可证副本替代
各文档分别包含的许可证副本，但条件是在所有其他方面，对每一份文档都遵守本许可证关于逐字
复制的规则。

您可以从此类合集中取出一份文档，并依本许可证单独分发，但必须在取出的文档中加入一份本许可
证副本，并在有关该文档逐字复制的所有其他方面遵守本许可证。


\section{与独立作品聚合（AGGREGATION WITH INDEPENDENT WORKS）}
\label{fdl-zh-section-aggregation-with-independent-works}


若将文档或其衍生作品与其他相互独立的文档或作品一同汇编在某一存储或分发媒介卷册中或其上，
且汇编产生的版权并未被用于在各独立作品所允许的范围之外限制汇编物使用者的法律权利，则该汇编
称为“聚合物”。当文档被纳入聚合物时，本许可证不适用于聚合物中本身并非文档衍生作品的其他作品。

若第 3 节的封面文字要求适用于这些文档副本，并且文档少于整个聚合物的一半，则文档的封面文字
可以放在聚合物内部包围该文档的封面上；若文档采用电子形式，则可放在相当于封面的电子界面上。
否则，封面文字必须出现在包围整个聚合物的印刷封面上。


\section{翻译（TRANSLATION）}
\label{fdl-zh-section-translation}


翻译视为一种修改，因此您可以依第 4 节的条款分发文档的译本。以译文替换不变章节，须取得
相应版权持有人的特别许可；但您可以在这些不变章节的原文之外，附加其中一部分或全部的译文。
您可以附上本许可证、文档中全部许可声明以及任何担保免责声明的译文，但同时必须包含本许可证
的英文原版，以及上述声明和免责声明的原版。若译文与本许可证、声明或免责声明的原版有任何
分歧，以原版为准。

若文档中有标题为“Acknowledgements”（致谢）、“Dedications”（献辞）或“History”（历史）
的章节，则第 4 节关于保留其标题（定义见第 1 节）的要求，通常意味着必须更改实际标题。


\section{终止（TERMINATION）}
\label{fdl-zh-section-termination}


除本许可证明确允许的情形外，您不得复制、修改、再许可或分发文档。任何其他复制、修改、
再许可或分发文档的企图均属无效，并将自动终止您依本许可证享有的权利。不过，只要从您处
获得副本或权利的各方继续完全遵守本许可证，其许可就不会因此终止。


\section{本许可证的未来修订版（FUTURE REVISIONS OF THIS LICENSE）}
\label{fdl-zh-section-future-revisions-of-this-license}


自由软件基金会可能不时发布 GNU 自由文档许可证的新修订版。新版本在精神上将与本版本相似，
但可能在细节上有所不同，以处理新的问题或关切。参见 https://www.gnu.org/copyleft/。

本许可证的每一版本均有不同的版本号。若文档说明本许可证的某个指定编号版本“或任何后续
版本”适用于该文档，您可以选择遵守该指定版本的条款和条件，也可以选择遵守自由软件基金会
此后正式发布（而非作为草案发布）的任何后续版本。若文档未指定本许可证的版本号，您可以
选择自由软件基金会曾经正式发布（而非作为草案发布）的任何版本。


\section{附录：如何将本许可证用于您的文档（ADDENDUM: How to use this License for your documents）}
\label{fdl-zh-section-addendum-how-to-use-this-license-for-your-documents}

若要将本许可证用于您撰写的文档，请在文档中包含一份本许可证副本，并在标题页之后紧接着
加入以下版权和许可声明：

\bigskip
\begin{quote}
    Copyright \copyright  YEAR  YOUR NAME.
    Permission is granted to copy, distribute and/or modify this document
    under the terms of the GNU Free Documentation License, Version 1.2
    or any later version published by the Free Software Foundation;
    with no Invariant Sections, no Front-Cover Texts, and no Back-Cover Texts.
    A copy of the license is included in the section entitled "GNU
    Free Documentation License".
\end{quote}
\bigskip

若文档包含不变章节、前封面文字和后封面文字，请用以下文字替换“with...Texts.”一行：

\bigskip
\begin{quote}
    with the Invariant Sections being LIST THEIR TITLES, with the
    Front-Cover Texts being LIST, and with the Back-Cover Texts being LIST.
\end{quote}
\bigskip

若文档有不变章节而没有封面文字，或者三者存在其他组合，请合并上述两种写法以符合实际情况。

若您的文档含有实质性的程序代码示例，我们建议同时依您选择的自由软件许可证（例如 GNU 通用
公共许可证）发布这些示例，使其可以用于自由软件。

%---------------------------------------------------------------------
